% !TeX root = ../../Book2.tex
% 纲要标题：## 第11章　Serre 谱序列与纤维化计算
% 纲要位置：几何学计划纲要.md，第 1429 行。
% 正文状态：初稿；尚未正式审核。

% 纲要内容开始
% 
% > 本章目标与证明范围：由基底骨架拉回滤过构造 Serre 谱序列，先解释局部系数再做计算。代表定理包括乘法、传递和所选有限生成结论；不扩充为所有谱序列的名录。[R1]
% 
% 纲要内容结束

\chapter{Serre 谱序列与纤维化计算}
\label{g2:ch:11}
\BookChapterNumber{g2:ch:11}

\begin{chapterabstract}
	基底的骨架把纤维化的总空间分成逐层可处理的部分。本章从这一几何滤过构造 Serre 谱序列，以沿路运输解释第二页的局部系数，并证明乘法、边缘映射、传递与比较。Hopf 丛、定向球丛、路径纤维化及酉群的计算将非零微分与环结构联系起来；纤维定向类的传递给出 Euler 类，回路球面的计算产生整数除幂代数。最后利用有限生成群在谱序列中的封闭性质，以及逐次消去高阶同伦群的有限截段，证明单连通 CW 复形的同调逐次有限生成当且仅当其同伦群逐次有限生成。
\end{chapterabstract}

\begin{learninggoals}
	\begin{itemize}
		\item 由基底骨架逆像构造几何滤过，识别相对群、运输作用及第二页，并说明固定总次数的收敛。
		\item 使用总次数的 Leibniz 规则和两条边缘映射，把页面上的微分与投影、限制及延拓障碍对应起来。
		\item 计算 Hopf 丛与球丛，定义 Euler 类并证明相应上同调正合列。
		\item 求出回路球面的加性群与所选杯积环，证明酉群、特殊酉群的整数外代数结构。
		\item 由谱序列证明有限生成性在纤维化中的传播及单连通空间的 Serre 有限生成定理。
	\end{itemize}
\end{learninggoals}

第8章的纤维化长正合列把 \(F\to E\to B\) 的同伦群联在一起。对同调，直接写出同样一条长正合列通常行不通。最简单的乘积已经说明原因：\(S^2\times S^1\) 的三次上同调来自一个二次类与一个一次类相乘，它同时使用基底和纤维的信息。Serre 谱序列先把这类混合信息放在坐标 \((2,1)\)，再检查总空间的粘合是否允许它保留下来。

\ProseParagraphBreak

比较两个纤维化
\[
	S^1\longrightarrow S^2\times S^1\longrightarrow S^2,
	\qquad S^1\longrightarrow S^3\longrightarrow S^2.
\]
它们的基底与纤维相同，因此第二页都有四个整数生成元；总空间的上同调却不同。第一张纤维化中，圆周类能够延拓到整个乘积；第二张中，它的延拓受阻，产生从 \((0,1)\) 到 \((2,0)\) 的非零微分。这是本章贯穿的例子：第二页描述局部材料，高阶微分记录这些材料在整体上拼接时遇到的关系。

\ProseParagraphBreak

本章主要沿用第10章的上同调指标 \(\mathrm d_r:(p,q)\mapsto(p+r,q-r+1)\)。有限生成证明另用同调版，其箭头改为 \((p,q)\mapsto(p-r,q+r-1)\)，并用 \(E^r_{p,q}\) 区分页面。局部系统沿用第5章“先走 \(\alpha\) 再走 \(\beta\)”的运输约定；第8章提供提升及路径纤维化，第9章提供 CW 逼近、Hurewicz 与消去同伦群的胞腔构造。关于谱序列及其几何应用的延伸阅读，可参看\cite[第三章，Spectral Sequences and Applications]{BottTu1982}。

% !TeX root = ../../Book2.tex
% 纲要标题：### 11.1 几何滤过及第二页
% 纲要位置：几何学计划纲要.md，第 1383 行。
% 正文状态：未撰写；本文件目前是框架。

% 纲要内容开始
% 
% * 基底骨架的逆像
% * 纤维局部平凡模型与相对群
% * 运输作用及局部系数
% * 第二页的识别与收敛条件
% 
% 纲要内容结束

\section{几何滤过及第二页}
\label{g2:sec:1101}

待撰写

\subsection{基底骨架的逆像}
\label{g2:subsec:1101:01}
待撰写

\subsection{纤维局部平凡模型与相对群}
\label{g2:subsec:1101:02}
待撰写

\subsection{运输作用及局部系数}
\label{g2:subsec:1101:03}
待撰写

\subsection{第二页的识别与收敛条件}
\label{g2:subsec:1101:04}
待撰写
% 纲要条目排为编号小节标题，并进入目录与书签；说明段仍保留于纲要注释中。

% !TeX root = ../../Book2.tex
% 纲要标题：### 11.2 乘积、传递与自然性
% 纲要位置：几何学计划纲要.md，第 1440 行。
% 正文状态：初稿；尚未正式审核。

% 纲要内容开始
% 
% * 杯积和 Leibniz 规则
% * 边缘映射的几何意义
% * 传递映射的代表元解释
% * 映射纤维化之间的比较
% 
% 纲要内容结束

\section{乘积、传递与自然性}
\label{g2:sec:1102}

第二页给出了许多群，真正的计算还需要它们之间的联系。乘法使一个生成元上的微分决定整片表格；边缘映射把表格接回纤维包含与基底投影；传递则把纤维上的一个类转为基底上的障碍。下面始终使用含幺交换系数环 \(R\)，谈到零行的普通系数时另假设纤维道路连通。

\subsection{杯积和 Leibniz 规则}
\label{g2:subsec:1102:01}

纤维的杯积与沿路运输交换，所以 \(\mathcal H^*(F;R)\) 是分次代数的局部系统。在局部平凡的记法中，\(a\otimes u\) 表示一个基底 \(p\) 次类配上一个纤维 \(q\) 次类。将两个这样的类相乘，必须先把第一个纤维因子越过第二个基底因子。因此第二页的乘法为
\begin{equation}\label{g2:eq:serre-product-sign}
	(a\otimes u)(b\otimes v)=(-1)^{q s}(a\smile b)\otimes(u\smile v),\qquad |a|=p,\ |u|=q,\ |b|=s.
\end{equation}
这是有张量识别时的公式。一般局部系数情形，则先按第5章的道路运输把系数移到同一纤维，再作纤维杯积；同一个 \((-1)^{qs}\) 仍须保留。

\begin{theorem}[Serre 谱序列的乘法]\label{g2:thm:serre-multiplicative}
	上同调 Serre 谱序列从第二页起是按总次数分次交换的代数谱序列。对 \(x\in E_r^{p,q}\)、\(y\in E_r^{s,t}\)，有
	\begin{equation}\label{g2:eq:serre-leibniz}
		\mathrm d_r(xy)=(\mathrm d_rx)y+(-1)^{p+q}x\,\mathrm d_ry.
	\end{equation}
	下一页的乘法由本页乘法诱导；第二页乘法由局部系数杯积和\cref{g2:eq:serre-product-sign}给出。目标上同调的滤过满足 \(F^pH^m(E)F^sH^n(E)\subset F^{p+s}H^{m+n}(E)\)，且极限页作为代数同构于这个环的伴随分次。
\end{theorem}
\begin{proof}
	先把几何滤过落实为一个可作乘法的链模型。用\cref{g2:lem:serre-filtered-cw-model}的相容 CW 逼近 \(X_p\to E_p\)，记 \(X=\bigcup_pX_p\)、\(C_*=C_*^{\mathrm{cell}}(X;\mathbb Z)\)。给每个胞腔记录它首次出现的层数；\(C_*(X_p)\) 即由层数不超过 \(p\) 的胞腔生成。积空间使用胞腔积的弱拓扑，积胞腔的层数是两因子层数之和。第4章的交叉积将其胞腔链识别为 \(C_*\otimes C_*\)，边界为 \(\partial(c\otimes d)=\partial c\otimes d+(-1)^{|c|}c\otimes\partial d\)。

	取基底对角映射的胞腔逼近 \(\Delta_c:B\to B\times B\)。它满足
	\[
		\Delta_c(B^k)\subset(B\times B)^k=\bigcup_{i+j\le k}B^i\times B^j.
	\]
	在 Hurewicz 替换中，将 \(\Delta_B\simeq\Delta_c\) 的同伦从总空间的真实对角映射开始提升，得到 \(\widetilde\Delta:E\to E\times E\)，且 \(\widetilde\Delta(E_k)\subset\bigcup_{i+j\le k}E_i\times E_j\)。在上述 CW 模型上逐层比较并作胞腔逼近，得到链映射
	\[
		D:C_*\longrightarrow C_*\otimes C_*,\qquad
		D(C_*(X_k))\subset\sum_{i+j\le k}C_*(X_i)\otimes C_*(X_j).
	\]
	这里逐层比较的依据是：相容 CW 逼近在每层都是弱同伦等价；对一个待延拓的胞腔，比较映射及比较同伦的障碍是该弱等价的一个相对同伦类，因而为零。按胞腔维数延拓，并在已处理子复形上固定选择，就得到相容比较。对有限个积因子的并也同样逐个积胞腔处理。这说明所用链对角确实保留总层数，而不只是普通上同调上的某个对角。

	令 \(C^*=\operatorname{Hom}(C_*,R)\)，滤过取在 \(C_*(X_{p-1})\) 上为零的上链。定义 \(uv=(u\otimes v)D\)。若 \(u\in F^p\)、\(v\in F^s\)，对层数小于 \(p+s\) 的任一积胞腔，至少有一个因子被相应上链消去，所以 \(uv\in F^{p+s}\)。\(D\) 为链映射又给出严格的上链恒等式
	\[
		\mathrm d(uv)=(\mathrm du)v+(-1)^{|u|}u\,\mathrm dv.
	\]
	若 \(u,v\) 是第 \(r\) 页的代表，即 \(\mathrm du\in F^{p+r}\)、\(\mathrm dv\in F^{s+r}\)，右端属于 \(F^{p+s+r}\)。更换代表时，乘法与上式使新增项落入第10章商公式的分母：多一层的近似闭元仍多一层，上边界乘积则改写为一个上边界加一个更高滤过项。因此乘积良定义，并且直接取上式在目标商中的类便得到\cref{g2:eq:serre-leibniz}。此论证对每个 \(r\) 适用，无需从第一微分的导子性质猜测后续微分。

	还要说明结合律与交换律为何在第二页成立。对角的两种三重迭代、交换两因子的对角，以及不同胞腔逼近，均由真实对角提供同伦。把同伦作相对胞腔逼近时，\(B^k\times I\) 只需进入总骨架 \(k+1\)。故所得链同伦提高层数至多一；对偶的上链同伦 \(K\) 降低滤过至多一。对于第 \(r\ge2\) 页的闭代表，总滤过为 \(P\) 时，误差形如 \(\mathrm dK(w)+K(\mathrm dw)\)。第一项是来自 \(F^{P-1}\subset F^{P-r+1}\) 的允许边界，第二项属于 \(F^{P+r-1}\subset F^{P+1}\)，两者在该页均为零。这证明从第二页起的结合律、单位律和总次数交换律。链对角代表真实对角，所以在总上同调上得到通常杯积。

	最后，在一个 \(p\) 胞腔与一个 \(s\) 胞腔上，用基底相对定向类乘纤维类识别第一页。交换中间的 \(q\) 次纤维因子和 \(s\) 次基底因子给出 \((-1)^{qs}\)，故第二页为所述乘法。对真闭代表同样取乘积，所得滤过和极限页乘积便是目标杯积的伴随分次。
\end{proof}

计算时需要区分两个次数：乘积交叉符号是 \((-1)^{qs}\)，导子公式的符号是 \((-1)^{p+q}\)。后者使用整个类的总次数。

\subsection{边缘映射的几何意义}
\label{g2:subsec:1102:02}

设 \(F\xrightarrow{i}E\xrightarrow{p}B\)，其中 \(B,F\) 道路连通。第二页零行是 \(H^*(B;R)\)，零列是纤维上同调中沿全部闭路运输不变的部分：
\[
	E_2^{n,0}=H^n(B;R),\qquad E_2^{0,n}=H^n(F;R)^{\pi_1(B)}.
\]
零行上的类只能被别的类击中，零列上的类只能向外发出微分。因此前者逐页取商，后者逐页取子群。

\begin{proposition}[边缘映射]\label{g2:prop:serre-edge-maps}
	Serre 谱序列中的两条边缘分别分解投影拉回与纤维限制：
	\[
		H^n(B;R)\twoheadrightarrow E_\infty^{n,0}=F^nH^n(E;R)\hookrightarrow H^n(E;R),
	\]
	\[
		H^n(E;R)\twoheadrightarrow E_\infty^{0,n}\hookrightarrow H^n(F;R)^{\pi_1(B)}\hookrightarrow H^n(F;R).
	\]
	两条复合依次为 \(p^*\) 和 \(i^*\)。特别地，能从总空间限制下来的纤维类必为永久圈。
\end{proposition}
\begin{proof}
	把原纤维化映到恒等纤维化 \(B\to B\)。后者仅有零行，谱序列在第二页退化，诱导的总上同调映射为 \(p^*\)；对骨架对的拉回以及相应正合偶取映射，即得第一条分解。由于总次数 \(n\) 中 \(F^{n+1}=0\)，最右的分级商就是实际子群 \(F^n\)。

	再从基点上的纤维化 \(F\to\{b_0\}\) 映入原纤维化。零骨架上的限制把全局类送到各纤维上的相容类，沿边的相容性正是运输不变性。滤过的定义给出 \(F^1H^n(E)=\ker i^*\)：若基底有多个顶点，沿路径运输表明在一个纤维为零等价于在全部顶点纤维为零。于是 \(H^n(E)/F^1\) 就是限制像，得到第二条分解。
\end{proof}

\begin{corollary}\label{g2:cor:serre-section-bottom-row}
	若 \(p\) 有截面，则所有落入零行的微分均为零。
\end{corollary}
\begin{proof}
	截面 \(s\) 满足 \(s^*p^*=\operatorname{id}\)，所以 \(p^*\) 单射。由第一条边缘分解，零行任何类都不能在逐页取商时被消去。
\end{proof}

这个判据常比直接追踪代表元快捷。它控制的是零行；其他行之间仍可能存在微分。

\subsection{传递映射的代表元解释}
\label{g2:subsec:1102:03}

设 \(q>0\)。一个纤维类先要保持运输不变，才能进入 \(E_2^{0,q}\)；随后还要通过 \(\mathrm d_2,\ldots,\mathrm d_q\) 的检验，才能到达第 \(q+1\) 页。最后可能发出的箭头恰为
\begin{equation}\label{g2:eq:serre-transgression}
	\tau=\mathrm d_{q+1}:E_{q+1}^{0,q}\longrightarrow E_{q+1}^{q+1,0}.
\end{equation}
它称为\term[chuandi]{传递}[transgression]。源是 \(H^q(F;R)\) 的一个子群，目标是 \(H^{q+1}(B;R)\) 的一个商；较早微分可能已在零行中加入关系。符号 \(\tau a\in H^{q+1}(B;R)\) 只有在该商确实等于整个基底群时才直接使用。

\begin{proposition}[传递的延拓解释]\label{g2:prop:serre-transgression}
	类 \(a\in H^q(F;R)^{\pi_1(B)}\) 属于 \(E_{q+1}^{0,q}\)，当且仅当它可延拓为 \(E_q=p^{-1}(B^q)\) 上的 \(q\) 次上同调类。选择这样的延拓后，它在 \((q+1)\) 胞腔上继续延拓的障碍组成一个基底上链；该上链所定义的类，模去较早微分的像，就是 \(\tau a\)。
\end{proposition}
\begin{proof}
	先对限制纤维化 \(E_q\to B^q\) 使用上一命题。零列 \((0,q)\) 的所有可能非零出射在第 \(q\) 页以前已经出现；其余箭头落在基底维数以外。因此这个位置的永久圈恰为 \(E_{q+1}^{0,q}\)，而边缘命题把它识别为 \(H^q(E_q;R)\to H^q(F;R)\) 的像。

	用一个闭上链 \(A\) 表示所选延拓，再将其作为上链延拓到 \(E_{q+1}\)。余边界 \(\mathrm d\widetilde A\) 在 \(E_q\) 上为零，故定义
	\[
		[\mathrm d\widetilde A]\in H^{q+1}(E_{q+1},E_q;R)
		\cong C^{q+1}(B;\mathcal H^0(F;R))=C^{q+1}(B;R).
	\]
	这个等式描述每个新胞腔上的障碍值。再取一次余边界为零，故这些值组成胞腔上圈。更换上链延拓只改动相对边界；更换较低骨架上的修正，则加入前面几页已经允许的边界。第10章代表元公式把这些不确定性恰好除去，并把剩余的 \(\mathrm d\widetilde A\) 识别为 \(\mathrm d_{q+1}[a]\)。
\end{proof}

还可用空间对表达一种十分实用的判定。记 \(\delta:H^q(F;R)\to H^{q+1}(E,F;R)\) 为连接同态。若 \(b\in H^{q+1}(B,b_0;R)\) 满足
\begin{equation}\label{g2:eq:transgression-relative-criterion}
	\delta a=p^*b,
\end{equation}
则 \(a\) 可传递，且 \(\tau a\) 是 \(b\) 在 \(E_{q+1}^{q+1,0}\) 中的像。事实上，在相容胞腔模型中取代表并调整一个相对上边界，就能使 \(\mathrm dA=p^*b\)；右端在基底 \(q\) 骨架上为零，于是 \(A\) 先给出上述延拓，再在下一层留下 \(b\)。连接同态采用“延拓后取余边界”的约定，因此这里没有额外负号。

\ProseParagraphBreak

例如，若 \(\widetilde H^j(F;R)=0\) 对 \(0<j<q\) 成立，则从 \((0,q)\) 出发的较早微分及落入 \((q+1,0)\) 的较早微分都没有非零源或靶。此时所有运输不变的 \(q\) 次类都可传递，目标也是整个 \(H^{q+1}(B;R)\)。球丛的特征类正是由这种情形产生。

\subsection{映射纤维化之间的比较}
\label{g2:subsec:1102:04}

给定交换图
\[
	\begin{tikzcd}
		F\arrow[r]\arrow[d,"f_F"]&E\arrow[r,"p"]\arrow[d,"f_E"]&B\arrow[d,"f_B"]\\
		F'\arrow[r]&E'\arrow[r,"p'"]&B',
	\end{tikzcd}
\]
沿路提升与 \(f_E\) 相容到同伦，因此纤维上同调的拉回给出局部系统映射 \(f_B^*\mathcal H^q(F';R)\to\mathcal H^q(F;R)\)。这才是第二页映射的完整输入；当基本群作用不平凡时，只列两个抽象群之间的同态还不够。

\begin{proposition}[自然性与比较]\label{g2:prop:serre-comparison}
	上述图诱导反向的谱序列映射 \(f_r^*:E_r'{}^{p,q}\to E_r^{p,q}\)，与微分、乘法及传递交换；第二页为局部系数上同调上的诱导映射，极限页为滤过映射 \(f_E^*\) 的伴随分次。如果第二页在所有位置都是同构，则 \(f_E^*:H^*(E';R)\to H^*(E;R)\) 是同构。
\end{proposition}
\begin{proof}
	先将 \(f_B\) 胞腔逼近，并在 Hurewicz 替换中提升这段同伦以相应调整 \(f_E\)。所得映射保持骨架逆像，故给出空间对正合列之间的交换图。第10章的导出正合偶构造于是逐页给出 \(f_r^*\)，且与微分交换。逐胞腔的相对群识别使用纤维运输，所以第一微分及第二页映射恰好是所述局部系数映射。对角映射的自然性使乘积相容；传递本身是一条指定微分，故也相容。

	若 \(f_2^*\) 为同构，取各页上同调归纳得到所有 \(f_r^*\) 为同构。固定总次数 \(n\)，两边的目标滤过只有 \(n+1\) 层。从最深一层开始，将相邻滤过项及其商放入短正合列，逐层用五引理，便得到 \(f_E^*\) 在 \(n\) 次的同构。这里使用的是前节已经证明的有限目标滤过。
\end{proof}

对同调版也有同样的比较，箭头方向与连续映射一致。特别地，若基底均单连通，\(f_B\) 与 \(f_F\) 都诱导整数同调同构，则 \(f_E\) 也诱导整数同调同构：第4章的同调万有系数短正合列先将基底同构提升为任意纤维同调系数下的同构，随后逐页比较即可。这个步骤解释了为什么“基底与纤维分别不变”能够推出总空间的同调不变。
% 纲要条目排为编号小节标题，并进入目录与书签；说明段仍保留于纲要注释中。

% !TeX root = ../../Book2.tex
% 纲要标题：### 11.3 球丛和路径纤维化
% 纲要位置：几何学计划纲要.md，第 1447 行。
% 正文状态：未撰写；本文件目前是框架。

% 纲要内容开始
% 
% * Hopf 丛中的非零微分
% * 球丛与 Euler 类的联系
% * 路径—回路纤维化
% * 悬挂和回路空间的所选计算
% 
% 纲要内容结束

\section{球丛和路径纤维化}
\label{g2:sec:1103}

待撰写

\subsection{Hopf 丛中的非零微分}
\label{g2:subsec:1103:01}
待撰写

\subsection{球丛与 Euler 类的联系}
\label{g2:subsec:1103:02}
待撰写

\subsection{路径—回路纤维化}
\label{g2:subsec:1103:03}
待撰写

\subsection{悬挂和回路空间的所选计算}
\label{g2:subsec:1103:04}
待撰写
% 纲要条目排为编号小节标题，并进入目录与书签；说明段仍保留于纲要注释中。

% !TeX root = ../../Book2.tex
% 纲要标题：### 11.4 经典群与分类空间
% 纲要位置：几何学计划纲要.md，第 1454 行。
% 正文状态：初稿；尚未正式审核。

% 纲要内容开始
% 
% * 酉群的整数上同调环
% * 通用圆周丛与 Euler 类
% * 维数、环结构及扭项的核验
% 
% 纲要内容结束

\section{经典群与分类空间}
\label{g2:sec:1104}

第八章把一个酉矩阵的第一列取出来，得到以低一阶酉群为纤维的球面丛。这张映射适合归纳：底球面只有两个非零上同调群，纤维的环结构又来自上一步。谱序列的乘法将使整页微分由少量生成元决定。随后考察总空间可缩的圆周丛，便能看到同一个纤维生成元怎样在另一种纤维化中传递成基底上的特征类。

\subsection{酉群的整数上同调环}
\label{g2:subsec:1104:01}

由\cref{g2:prop:classical-sphere-bundles}，取第一列给出局部平凡丛
\[
	\operatorname{U}(n-1)\xrightarrow{i_n}\operatorname{U}(n)\xrightarrow{p_n}S^{2n-1},
	\qquad i_n(A)=\operatorname{diag}(1,A).
\]
当 \(n\ge2\) 时，底球面单连通，所以纤维上同调的局部系统为常值。这里完全使用整数系数；保留整数才能同时判断生成元的倍数以及挠项是否出现。

\begin{theorem}[酉群的整数上同调环]\label{g2:thm:unitary-cohomology}
	对每个 \(n\ge1\)，可选齐次类 \(x_{2j-1}\in H^{2j-1}(\operatorname{U}(n);\mathbb Z)\)（\(1\le j\le n\)），使
	\[
		H^*(\operatorname{U}(n);\mathbb Z)\cong\Lambda_{\mathbb Z}(x_1,x_3,\ldots,x_{2n-1}).
	\]
	这些选择可以逐个 \(n\) 相容，使 \(i_n^*x_{2j-1}\) 为低一阶酉群中的同名类（\(j<n\)），而 \(x_{2n-1}=p_n^*\omega_{2n-1}\)，其中 \(\omega_{2n-1}\) 是球面的一个整数上同调生成元。这里的外代数具有关系 \(x_ix_j=-x_jx_i\)、\(x_i^2=0\)；每个生成元至多出现一次的递增单项式构成其整数基。
\end{theorem}
\begin{proof}
	初始情形为 \(\operatorname{U}(1)=S^1\)，其环是 \(\Lambda_{\mathbb Z}(x_1)\)。设已知 \(\operatorname{U}(n-1)\) 的结论。记 \(m=2n-1\)，并将纤维中的生成元暂记为 \(\bar x_1,\ldots,\bar x_{m-2}\)。基底球面的同调自由，故第二页及其乘法为
	\[
		E_2^{*,*}\cong H^*(S^m;\mathbb Z)\otimes\Lambda_{\mathbb Z}(\bar x_1,\ldots,\bar x_{m-2}).
	\]
	只有第 \(0,m\) 两列可能非零。由 \(\mathrm d_r\) 的双次数，唯一可能非零的微分为 \(\mathrm d_m\)，从第零列指向第 \(m\) 列。对每个旧生成元，
	\[
		\mathrm d_m:E_m^{0,2j-1}\longrightarrow E_m^{m,2j-m},\qquad j<n,
	\]
	目标的第二指标至多为 \(-1\)，所以 \(\mathrm d_m\bar x_{2j-1}=0\)。底球面的类也没有非零出微分。\cref{g2:thm:serre-multiplicative}的 Leibniz 规则于是说明，所有生成元的乘积都具有零微分。这些乘积生成整页，故 \(E_2=E_\infty\)。注意，纤维本身可能在很高次数仍有非零上同调；真正使微分全部消失的是低次生成元与导数性共同作用。

	每个 \(E_\infty^{p,q}\) 都是有限秩自由阿贝尔群。\cref{g2:prop:free-graded-splitting}说明 \(H^k(\operatorname{U}(n);\mathbb Z)\) 也自由：其有限滤过的每个商都有可提升的基。边缘映射的描述则给出满射
	\[
		i_n^*:H^*(\operatorname{U}(n);\mathbb Z)\longrightarrow H^*(\operatorname{U}(n-1);\mathbb Z).
	\]
	据此选择 \(x_{2j-1}\) 提升 \(\bar x_{2j-1}\)。在另一条边上，\(p_n^*\omega_m\) 的伴随分次类就是 \(E_\infty^{m,0}\) 的生成元，令它为 \(x_m\)。由于所有 \(x_{2j-1}\) 都是奇次类，杯积的分次交换律给出 \(2x_{2j-1}^2=0\)；总上同调无挠，故每个平方为零。这就构造出从所述外代数到 \(H^*(\operatorname{U}(n);\mathbb Z)\) 的环同态。

	最后核对这些乘积确实构成基，以解决乘法延拓。在固定总次数 \(k\) 上，滤过只有第零列和第 \(m\) 列两个可能非零的商。不含 \(x_m\) 的递增单项式限制到纤维，正好给出第零列的一组基；含 \(x_m\) 的单项式可写成 \(x_m x_I\)，其在第 \(m\) 列的类为 \(\omega_m\bar x_I\)，也正好给出该列的一组基。任意上同调类先减去第一组基的线性组合，就落在第 \(m\) 层，再由第二组表示；任何线性关系先限制到纤维，第一组的系数都为零，再投到第 \(m\) 层，第二组的系数也都为零。因此这张环同态为同构。沿归纳过程选择提升，即得所述相容性。
\end{proof}

这个证明给出的生成元有两种来源：旧生成元从纤维提升，新生成元从底球面拉回。例如 \(\operatorname{U}(2)\) 的基为 \(1,x_1,x_3,x_1x_3\)，分别处于次数 \(0,1,3,4\)。其中 \(x_3=p_2^*\omega_3\)，而 \(x_1\) 可取行列式映射 \(\det:\operatorname{U}(2)\rightarrow S^1\) 拉回的一次生成元；行列式在标准 \(\operatorname{U}(1)\) 纤维上恰为恒等映射。

\begin{corollary}[特殊酉群的整数上同调环]\label{g2:cor:special-unitary-cohomology}
	对 \(n\ge2\)，有
	\[
		H^*(\operatorname{SU}(n);\mathbb Z)\cong\Lambda_{\mathbb Z}(y_3,y_5,\ldots,y_{2n-1}),\qquad |y_j|=j.
	\]
\end{corollary}
\begin{proof}
	\cref{g2:prop:classical-sphere-bundles}给出 \(\operatorname{SU}(2)\cong S^3\) 以及 \(\operatorname{SU}(n-1)\rightarrow\operatorname{SU}(n)\rightarrow S^{2n-1}\)。从 \(n=3\) 开始归纳。纤维的生成元处于次数 \(3,5,\ldots,2n-3\)，所以唯一可能的 \(\mathrm d_{2n-1}\) 在每个生成元上都有负的目标第二指标，因而全为零；导数性使全页退化。全部分次块自由，纤维生成元可以提升，新生成元由球面拉回。无挠性使奇次平方为零，含新生成元及不含新生成元的单项式分别给出两层的基。由上一证明中的逐层基论证，得到所述外代数。
\end{proof}

\subsection{通用圆周丛与 Euler 类}
\label{g2:subsec:1104:02}

在上面的酉群计算中，旧生成元全部存活。现在把圆周放进总空间可缩的纤维化，其一次类就必须被消去。消去它的微分落在基底二次上同调中；自然性将把这个类送到每个具体圆周丛的基底上。

\ProseParagraphBreak

取有限支撑复序列构成的单位球面
\[
	S^\infty=\bigcup_{N\ge1}S^{2N-1}\subset\mathbb C^\infty,
\]
赋予与有限维球面相容的 CW 弱拓扑。这里可从实球面的逐维赤道与半球胞腔结构取并，每个有限子复形都包含于足够大的有限维球面。标量右乘给出自由 \(S^1\) 作用，其商为 \(\mathbb CP^\infty\)。在 \(z_j\ne0\) 的坐标片上，仍可选择第 \(j\) 坐标为正实数的唯一单位代表；有限阶段的这些公式相容，给出局部平凡丛
\begin{equation}\label{g2:eq:universal-circle-bundle}
	S^1\longrightarrow S^\infty\xrightarrow{\pi}\mathbb CP^\infty.
\end{equation}
其中 \(S^\infty\) 可缩。事实上，由\cref{g2:thm:finite-subcomplex}，任一球面映射的像包含在有限阶段；再放进维数大于定义域维数的球面，由\cref{g2:prop:sphere-connectivity}可知它零伦。因此 \(S^\infty\) 道路连通且全部正阶同伦群为零，\cref{g2:thm:whitehead}用于点到这个 CW 复形的映射，便得到可缩性。

\ProseParagraphBreak

基底上的\term[tautological-complex-line-bundle]{典范复线丛}[tautological complex line bundle]为
\[
	\gamma^1=\{(\ell,v):\ell\in\mathbb CP^\infty,\ v\in\ell\}\longrightarrow\mathbb CP^\infty.
\]
它的单位圆周丛就是\eqref{g2:eq:universal-circle-bundle}。复线的定向使其底层实二维丛 \(\gamma^1_{\mathbb R}\) 定向；按\cref{g2:def:spherical-euler-class}，记
\[
	u\coloneqq e(\gamma^1_{\mathbb R})=\mathrm d_2(a),\qquad a\in H^1(S^1;\mathbb Z)
\]
为正向纤维生成元的传递。这个 \(u\) 是 \(H^2(\mathbb CP^\infty;\mathbb Z)\cong\mathbb Z\) 的生成元：\(E_2^{0,1}\) 与 \(E_2^{2,0}\) 之间只有 \(\mathrm d_2\)，而总空间的一次、二次上同调都为零，迫使该微分既单射又满射。它与\cref{g2:thm:complex-projective-ring}选用的二次生成元至多相差符号。

\ProseParagraphBreak

由该定理及乘法结构，第二页和微分完整写成
\[
	E_2^{*,*}=\mathbb Z[u]\otimes\Lambda_{\mathbb Z}(a),\qquad |u|=(2,0),\quad |a|=(0,1),
\]
\[
	\mathrm d_2(u^k)=0,\qquad \mathrm d_2(u^ka)=u^{k+1}\quad(k\ge0).
\]
每个正次数的底类都被上面一行的类消去，只留下单位元，正好对应总空间可缩。在有限维 Hopf 丛中，底环被 \(u^{N}=0\) 截断，因而 \(u^{N-1}a\) 留下总空间 \(S^{2N-1}\) 的顶次类；通用丛把这些有限截断相容地接在一起。

\begin{proposition}[复线丛的拉回模型]\label{g2:prop:line-bundle-pullback-model}
	设 \(B\) 为有限 CW 复形，\(L\rightarrow B\) 为连续复线丛。则存在某个 \(N\) 及连续映射 \(f:B\rightarrow\mathbb CP^{N-1}\subset\mathbb CP^\infty\)，使 \(L\cong f^*\gamma^1\)。把 \(L\) 看成定向实二维丛后，其 Euler 类满足
	\[
		e(L_{\mathbb R})=f^*u.
	\]
\end{proposition}
\begin{proof}
	有限 CW 复形紧且 Hausdorff，故可取有限平凡化覆盖 \(U_1,\ldots,U_N\) 及从属的连续单位分解 \(\rho_1,\ldots,\rho_N\)。这里所需的事实是 \(\rho_j\ge0\)、\(\operatorname{supp}\rho_j\subset U_j\) 及 \(\sum_j\rho_j=1\)；存在性见第一卷第4.4节的连续单位分解定理。局部 Hermitian 内积也可用这些权重相加，得到 \(L\) 上的连续正定内积，因为非零向量的各项平方长度都为正。将局部标架归一化，得到单位截面 \(s_j\)。

	对 \(b\in U_j\) 及单位向量 \(v\in L_b\)，唯一写成 \(v=s_j(b)\lambda_j(v)\)，其中 \(\lambda_j(v)\in S^1\)。令
	\[
		J(v)=\bigl(\sqrt{\rho_1(b)}\lambda_1(v),\ldots,\sqrt{\rho_N(b)}\lambda_N(v)\bigr)\in S^{2N-1},
	\]
	并在 \(b\notin U_j\) 时把相应坐标规定为零。支集条件使这些补零后的坐标连续，而 \(\sum\rho_j=1\) 保证所得向量确实为单位向量。对 \(z\in S^1\)，有 \(J(vz)=J(v)z\)，故 \(f(b)=[J(v)]\) 与同一纤维中的单位向量选择无关，并由局部截面看出 \(f\) 连续。

	将 \(J\) 在每条复线上复线性延拓，得到 \(L_b\) 到 \(\mathbb C^N\) 的等距单射，其像为 \(f(b)\) 所代表的复直线。因此
	\[
		L\longrightarrow f^*\gamma^1,\qquad v\longmapsto(b,J(v))
	\]
	是纤维上的复线性同构；在局部非零标架上，其逆由一个连续非零系数的倒数给出，故为丛同构。限制到单位向量便得到圆周丛的拉回交换图
	\[
		\begin{tikzcd}[column sep=large]
			S(L)\arrow[r,"J"]\arrow[d] & S^\infty\arrow[d,"\pi"]\\
			B\arrow[r,"f"'] & \mathbb CP^\infty.
		\end{tikzcd}
	\]
	纤维映射保持复定向。由\cref{g2:prop:spherical-euler-naturality}，其 Euler 类为 \(f^*u\)。
\end{proof}

这个构造解释了“通用”的含义：一个具体线丛的局部标架经加权后，确定一张到射影空间的映射，再由同一个典范线丛拉回。赋予每个丛一个随拉回自然变化的上同调类，正是\term[tezheng-lei]{特征类}[characteristic class]的基本形式。这里得到的 \(e(L_{\mathbb R})\) 已是这样一个类；它由圆周的一个整数生成元及传递构造而来。不同单位分解所产生的映射都给出同一个 Euler 类，因为它们拉回得到的是同一个丛。

\ProseParagraphBreak

在一般秩中，射影空间将由 Grassmann 空间替代，单位向量将由酉标架替代。\chapref{g2:ch:19}将建立相应分类空间与通用丛，\chapref{g2:ch:21}再系统构造复向量丛的特征类。眼前这张圆周丛已显示出关键计算方式：先在通用模型中求纤维类的传递，再利用拉回自然性把答案送到给定基底。

\subsection{维数、环结构及扭项的核验}
\label{g2:subsec:1104:03}

得到一个环的候选答案后，可以分别检验其次数、乘法及整数信息。对 \(\operatorname{U}(n)\)，上一小节的论证已经解决了后两者；把全部基元的次数汇总，还会产生几项直观的核验。

\begin{corollary}[酉群的 Betti 数与顶次类]\label{g2:cor:unitary-betti-polynomial}
	记 \(b_k\) 为整数上同调的秩。对 \(n\ge1\)，有
	\[
		\sum_k b_k(\operatorname{U}(n))t^k=\prod_{j=1}^n(1+t^{2j-1}),
		\qquad \sum_k b_k(\operatorname{U}(n))=2^n,
	\]
	且最高非零次数为 \(n^2\)，由 \(x_1x_3\cdots x_{2n-1}\) 生成，Euler 示性数为零。对 \(n\ge2\)，特殊酉群的相应多项式为 \(\prod_{j=2}^n(1+t^{2j-1})\)，最高非零次数为 \(n^2-1\)，Euler 示性数也为零。
\end{corollary}
\begin{proof}
	外代数的一组基由每个生成元“选取或不选取”得到，所以次数生成函数为所列乘积；令 \(t=1\) 得到总秩。最高次数是
	\[
		1+3+\cdots+(2n-1)=n^2,
	\]
	且只有选取全部生成元的一个单项式达到该次数。令 \(t=-1\)，每个奇次因子都为零，得到 Euler 示性数。特殊酉群删去一次生成元，故顶次降为 \(n^2-1\)，其余计算相同。
\end{proof}

顶次与矩阵的实参数数目一致。单位酉矩阵附近的无穷小变动由斜 Hermitian 矩阵给出：对角线有 \(n\) 个纯虚参数，每对非对角位置由一个复数决定，共有 \(n(n-1)\) 个实参数，合为 \(n^2\)。行列式为一增加迹为零的一个实条件，得到 \(n^2-1\)。这与上同调中唯一顶次单项式的位置吻合。

\ProseParagraphBreak

例如
\[
	\sum_kb_k(\operatorname{U}(3))t^k=(1+t)(1+t^3)(1+t^5)
	=1+t+t^3+t^4+t^5+t^6+t^8+t^9.
\]
次数为 \(4\) 的类是 \(x_1x_3\)，次数为 \(6\) 的类是 \(x_1x_5\)，次数为 \(8\) 的类是 \(x_3x_5\)。把这些乘积列出来，有助于避免将“只有奇次生成元”误读为“只有奇次上同调”。

\ProseParagraphBreak

对任一素数 \(\ell\)，系数短正合列 \(0\rightarrow\mathbb Z\xrightarrow{\times\ell}\mathbb Z\rightarrow\mathbb F_\ell\rightarrow0\) 给出
\[
	H^k(\operatorname{U}(n);\mathbb Z)\xrightarrow{\times\ell}H^k(\operatorname{U}(n);\mathbb Z)
	\longrightarrow H^k(\operatorname{U}(n);\mathbb F_\ell)
	\longrightarrow H^{k+1}(\operatorname{U}(n);\mathbb Z)\xrightarrow{\times\ell}H^{k+1}(\operatorname{U}(n);\mathbb Z).
\]
末箭头单射，故约化模 \(\ell\) 给出整数上同调张量 \(\mathbb F_\ell\) 到模 \(\ell\) 上同调的同构。杯积与系数映射相容，因而
\[
	H^*(\operatorname{U}(n);\mathbb F_\ell)\cong\Lambda_{\mathbb F_\ell}(\bar x_1,\bar x_3,\ldots,\bar x_{2n-1}).
\]
当 \(\ell=2\) 时，奇次平方为零的理由是这些类来自整数类的约化；单凭特征二下的分次交换律不能得到这个结论。

\begin{example}[旋转群中的整数挠类]\label{g2:exm:rotation-torsion-comparison}
	\cref{g2:lem:quaternion-rotation-cover}把 \(\operatorname{SO}(3)\) 识别为 \(S^3/\{\pm1\}=\mathbb RP^3\)。\cref{g2:prop:projective-integer-homology}给出
	\[
		H_0\cong H_3\cong\mathbb Z,\qquad H_1\cong\mathbb Z/2,\qquad H_2=0.
	\]
	整数万有系数定理据此得到
	\[
		H^k(\operatorname{SO}(3);\mathbb Z)\cong
		\begin{cases}
			\mathbb Z,&k=0,3,\\
			\mathbb Z/2,&k=2,\\
			0,&\text{其他次数}.
		\end{cases}
	\]
	其中二次的 \(\mathbb Z/2\) 是 \(\operatorname{Ext}_{\mathbb Z}^1(H_1,\mathbb Z)\)。有理系数却只剩 \(H^*(\operatorname{SO}(3);\mathbb Q)\cong\Lambda_{\mathbb Q}(z_3)\)。所以一个有理外代数答案仍可能遗漏整数挠项；酉群的整数外代数结论，依靠的是整数谱序列中自由分次块与生成元提升的完整论证。
\end{example}

% 纲要条目排为编号小节标题，并进入目录与书签；说明段仍保留于纲要注释中。

% !TeX root = ../../Book2.tex
% 纲要标题：### 11.5 有限生成与进一步方法
% 纲要位置：几何学计划纲要.md，第 1460 行。
% 正文状态：初稿；尚未正式审核。

% 纲要内容开始
% 
% * Serre 类与有限生成群
% * 同调与同伦群的有限生成性
% * Adams 与 Eilenberg–Moore 谱序列
% 
% ---
% 
% 纲要内容结束

\section{有限生成与进一步方法}
\label{g2:sec:1105}

谱序列也能回答比“这个群到底是什么”更粗、却更普遍的问题。例如，有限 CW 复形的同调由有限矩阵计算，但它的高阶同伦群并没有同样直接的矩阵表示。能否至少证明每个同伦群只需有限多个生成元？这一节用谱序列把有限生成性在纤维、基底与总空间之间传递，再用第9章的胞腔附着将同伦群逐层分离出来。以下未注明的同调均取整数系数。

\subsection{Serre 类与有限生成群}
\label{g2:subsec:1105:01}

一个页面上的群经取核、取像和取商，成为下一页的群；目标群又是有限个极限页项的逐次扩张。因此，适合被谱序列追踪的性质，应当在这几种操作下保持。

\begin{definition}[Serre 类]\label{g2:def:serre-class}
	阿贝尔群的一类 \(\mathcal C\)，若包含零群，且对任意短正合列 \(0\to A\to B\to C\to0\) 都有
	\[
		B\in\mathcal C\quad\Longleftrightarrow\quad A,C\in\mathcal C,
	\]
	就称为一个\term[Serre lei]{Serre 类}[Serre class]。这等价于对取子群、取商及有限次扩张封闭。若一个同态的核与余核都属于 \(\mathcal C\)，便称它为模 \(\mathcal C\) 的同构。
\end{definition}

有限生成阿贝尔群组成一个 Serre 类，记为 \(\mathcal C_{\mathrm{fg}}\)。其子群有限生成可由整数自由模的子模定理及有限生成群分解推出，商群显然由生成元的像生成。对扩张，提升商群的一组有限生成元，再加入子群的一组有限生成元，就得到中间群的有限生成集。有限阿贝尔群也组成一个 Serre 类：短正合列中中间群的阶是两端群阶的乘积。

\ProseParagraphBreak

Serre 类的定义本身没有规定张量积与 Tor 的行为。使用第二页时，还须核对这些系数运算。对 \(\mathcal C_{\mathrm{fg}}\)，将两群写成有限个无限循环群与有限循环群的直和，再用\cref{g2:prop:cyclic-ext-tor}，可知它们的张量积和 \(\operatorname{Tor}_1^{\mathbb Z}\) 仍有限生成。于是当 \(H_i(B)\) 与 \(A\) 都逐次有限生成时，同调万有系数短正合列
\begin{equation}\label{g2:eq:finite-generation-coefficients}
	0\longrightarrow H_p(B)\otimes A\longrightarrow H_p(B;A)
	\longrightarrow\operatorname{Tor}_1^{\mathbb Z}(H_{p-1}(B),A)\longrightarrow0
\end{equation}
说明系数同调也有限生成。下面用这一项具体性质，而不对任意 Serre 类作未经核对的推广。

\subsection{同调与同伦群的有限生成性}
\label{g2:subsec:1105:02}

称一个道路连通空间具有\term[tongdiao-youxianxing]{逐次有限生成同调}[degreewise finitely generated homology]，如果每个 \(H_n(-;\mathbb Z)\) 都有限生成。各次数可以有不同数量的生成元，也允许无限多个次数非零。譬如 \(\mathbb CP^\infty\) 满足此条件，尽管它有无限多个胞腔。

\begin{lemma}[纤维化中的有限生成性]\label{g2:lem:fibration-finite-generation}
	设 \(F\to E\to B\) 是道路连通空间的纤维化，基底为单连通 CW 复形。若 \(F,E,B\) 中任意两个具有逐次有限生成同调，则第三个也有。
\end{lemma}
\begin{proof}
	用同调谱序列
	\[
		E^2_{p,q}=H_p(B;H_q(F))\Longrightarrow H_{p+q}(E),\qquad
		\mathrm d_r:E^r_{p,q}\to E^r_{p-r,q+r-1}.
	\]
	单连通性消去了运输作用，且 \(H_0(F)=H_0(B)=\mathbb Z\)。每个固定位置经过有限多页便稳定，每个目标次数也只有有限层。

	先设 \(B,F\) 的同调逐次有限生成。式\eqref{g2:eq:finite-generation-coefficients}使全部第二页项有限生成。取核、取商保留此性质，故全部极限页项也有限生成；目标群由有限个这样的群扩张而成，结论随之成立。

	再设 \(F,E\) 满足条件，归纳证明 \(H_k(B)\) 有限生成。假定次数小于 \(k\) 的基底群已有此性质，则第二页中所有 \(p<k\) 的项都有限生成。同调零行没有来箭头，所以对 \(2\le r\le k\) 有
	\[
		0\longrightarrow E^{r+1}_{k,0}\longrightarrow E^r_{k,0}
		\xrightarrow{\mathrm d_r}\operatorname{im}\mathrm d_r\longrightarrow0,
	\]
	其中像是 \(E^r_{k-r,r-1}\) 的子群，故有限生成。稳定项 \(E^\infty_{k,0}\) 是 \(H_k(E)\) 的一个分级商，也有限生成。倒着经过有限多页，短正合列遂说明 \(E^2_{k,0}=H_k(B)\) 有限生成。\(k=1\) 时没有这些微分，结论直接来自稳定项；这也给出归纳起点。

	最后设 \(B,E\) 满足条件，归纳证明 \(H_k(F)\) 有限生成。零列没有去箭头，可能进入 \((0,k)\) 的微分只有
	\[
		\mathrm d_r:E^r_{r,k-r+1}\longrightarrow E^r_{0,k},\qquad 2\le r\le k+1.
	\]
	所有源项的纤维次数都小于 \(k\)，故按归纳假设及系数公式有限生成。各页有短正合列
	\[
		0\longrightarrow\operatorname{im}\mathrm d_r\longrightarrow E^r_{0,k}
		\longrightarrow E^{r+1}_{0,k}\longrightarrow0.
	\]
	从 \(H_k(E)\) 的有限生成分级商 \(E^\infty_{0,k}\) 倒推，便得 \(E^2_{0,k}=H_k(F)\) 有限生成。零次群为 \(\mathbb Z\)，提供了这里的归纳起点。
\end{proof}

\begin{corollary}\label{g2:cor:loop-finite-generation}
	若单连通 CW 复形 \(B\) 具有逐次有限生成同调，则 \(\Omega B\) 也有。
\end{corollary}
\begin{proof}
	将引理的第三种情形用于 \(\Omega B\to P_{b_0}B\to B\)。道路空间可缩，回路空间因基底单连通而道路连通。
\end{proof}

要把结论送到同伦群，需要先处理只含一个同伦群的空间。下述证明把所需模型与有限性理由一并写出。

\begin{lemma}[Eilenberg–MacLane 空间的有限生成性]\label{g2:lem:em-space-finite-generation}
	若 \(A\) 是有限生成阿贝尔群，则对每个 \(n\ge1\)，空间 \(K(A,n)\) 具有逐次有限生成整数同调。
\end{lemma}
\begin{proof}
	先说明这里选择哪个 CW 模型不影响结论。对 \(n\ge2\)，用第9章的稀疏 CW 模型把 \(K(A,n)\) 换成除一个顶点外没有低于 \(n\) 维胞腔的模型。其 \(n\) 骨架是球面楔和，且 \(H_n=C_n/\operatorname{im}\partial_{n+1}\)。给定另一个 \(K(A,n)\)，把每个 \(n\) 胞腔的球面送到它所代表的 \(A\) 中元素；Hurewicz 同构保证这样的代表存在。每个 \((n+1)\) 胞腔的边界在 \(A\) 中为零，故可延拓；更高胞腔的延拓障碍落在目标已消失的高阶同伦群中，也全为零。所得映射在唯一非零同伦群上为同构，Whitehead 定理给出同伦等价。对 \(n=1\)，先在一骨架实现基本群的指定同态，二胞腔的关系保证延拓，随后用目标高阶同伦群全零延拓。因而同样得到模型唯一性。非 CW 纤维可先取 CW 逼近；\cref{g2:lem:weak-equivalence-homology}保证所用同调不变。

	现处理 \(n=1\)。将 \(A\) 写成 \(\mathbb Z^r\times G\)，其中 \(G\) 是有限阿贝尔群。环面 \(T^r\) 是 \(K(\mathbb Z^r,1)\)，已有有限 CW 模型。对有限群 \(G\)，构造如下空间：在次数 \(k\) 为每个序列 \((g_0,\ldots,g_k)\) 放置一个 \(k\) 单形，按删去一个顶点的面映射粘合，并把相邻重复顶点的退化单形压到相应低维单形。所得实现记为 \(EG\)。它是 CW 复形，\(G\) 通过同时左乘全部顶点自由作用；商 \(BG\) 每维只有有限多个胞腔，因为每维的序列总数有限。

	\(EG\) 可缩有一个直接的公式。对由 \((g_0,\ldots,g_k)\) 与重心坐标 \((t_0,\ldots,t_k)\) 表示的点，送到由
	\[
		(1,g_0,\ldots,g_k),\qquad (s,(1-s)t_0,\ldots,(1-s)t_k)
	\]
	表示的点，\(0\le s\le1\)。它与面、退化的识别相容；\(s=0\) 是原点，\(s=1\) 是顶点 \(1\)。有限群的自由作用在 Hausdorff 空间上给出覆盖：对一点及其有限多个不同的平移，先取互不相交的小邻域，再有限次相交得到与所有非平凡平移不交的邻域。故 \(EG\to BG\) 是具有可缩总空间的通用覆盖，\(BG\) 是 \(K(G,1)\)。其每维胞腔链为有限秩自由群，所以每个同调群有限生成。乘积 \(T^r\times BG\) 的同伦群逐坐标计算，说明它是 \(K(A,1)\)；Künneth 定理及有限生成群的张量、Tor 性质给出其同调有限生成。

	对 \(n\ge2\) 作归纳。路径纤维化
	\[
		\Omega K(A,n)\longrightarrow P K(A,n)\longrightarrow K(A,n)
	\]
	的纤维由同伦长正合列具有 \(K(A,n-1)\) 的同伦群，取 CW 逼近并用刚证的模型唯一性，便知其同调逐次有限生成。总空间可缩，基底单连通。\cref{g2:lem:fibration-finite-generation}中“纤维与总空间推出基底”的情形完成归纳。
\end{proof}

接着把一个空间的同伦群分成有限层。这里先构造有限生成证明所用的部分；第12章将进一步研究层与层之间的附着信息。

\begin{lemma}[有限同伦截段]\label{g2:lem:finite-homotopy-truncations}
	设 \(X\) 为单连通带基点 CW 复形。对每个 \(n\ge1\)，存在包含 \(X\hookrightarrow X[n]\)，使
	\[
		\pi_i(X[n])\cong\pi_i(X)\ (i\le n),\qquad \pi_i(X[n])=0\ (i>n),
	\]
	且包含在 \(H_i\) 上对 \(i\le n\) 为同构。可以选择延拓恒等映射的映射 \(X[n]\to X[n-1]\)（\(n\ge2\)），其同伦纤维为 \(K(\pi_n(X),n)\) 的弱同伦型；\(X[1]\) 可缩。
\end{lemma}
\begin{proof}
	从 \(X\) 出发，依次用\cref{g2:lem:kill-homotopy-group}消去 \(\pi_{n+1},\pi_{n+2},\ldots\)。所加胞腔维数至少为 \(n+2\)，因此所有 \(i\le n\) 次同伦群保持不变，相对胞腔链在次数小于 \(n+2\) 时为零，故相应同调也保持。对无限并，第9章的紧像有限子复形论证保证固定次数的同伦群达到稳定值。\(n=1\) 时，初始基本群已为零，所得 \(X[1]\) 的全部同伦群均为零，Whitehead 定理使它可缩。

	在 \(X[n]\) 的子复形 \(X\) 上取包含到 \(X[n-1]\) 的映射。一个新增 \(d\) 胞腔满足 \(d\ge n+2\)，延拓所需消去的边界类属于 \(\pi_{d-1}(X[n-1])=0\)，所以可逐胞腔延拓到整个 \(X[n]\)。将所得映射作第8章的路径纤维化替换，记纤维为 \(F_n\)。在次数 \(i<n\) 上，映射诱导由 \(X\) 给出的同构；基底在次数 \(i\ge n\) 的群均为零。长正合列因而给出
	\[
		\pi_i(F_n)=0\ (i\ne n),\qquad \pi_n(F_n)\cong\pi_n(X).
	\]
	两空间单连通还保证纤维道路连通。取纤维的 CW 逼近，即得到指定的 Eilenberg–MacLane 同伦型。
\end{proof}

\begin{theorem}[Serre 有限生成定理]\label{g2:thm:serre-finite-generation}
	对单连通 CW 复形 \(X\)，下列两个条件等价：
	\begin{enumerate}
		\item 每个 \(H_n(X;\mathbb Z)\) 都有限生成；
		\item 每个 \(\pi_n(X)\)（\(n\ge2\)）都有限生成。
	\end{enumerate}
	特别地，单连通有限 CW 复形的全部同伦群都有限生成。
\end{theorem}
\begin{proof}
	先假定全部同伦群有限生成。\(X[1]\) 可缩，从它开始沿
	\[
		K(\pi_j(X),j)\longrightarrow X[j]\longrightarrow X[j-1]
	\]
	逐层归纳。这里的箭头按上一引理作纤维化替换，仍用原名表示同伦等价的总空间。Eilenberg–MacLane 引理使每一纤维的同调逐次有限生成，纤维化引理再把此性质传到 \(X[j]\)。固定 \(n\)，取 \(j\ge n\)，包含给出 \(H_n(X)\cong H_n(X[j])\)，所以 \(H_n(X)\) 有限生成。相同论证还说明：若只知道 \(\pi_i(X)\) 在 \(i\le j\) 有限生成，则 \(X[j]\) 的全部同调已经逐次有限生成。

	反过来，假定全部 \(H_i(X)\) 有限生成，归纳证明 \(\pi_n(X)\) 有限生成。起点 \(n=2\) 是 Hurewicz 同构 \(\pi_2(X)\cong H_2(X)\)。设低于 \(n\) 的同伦群已知有限生成。上一段的有限截段结论保证基底 \(Y=X[n-1]\) 的全部同调有限生成。令 \(Z=X[n]\)，考虑纤维化 \(F_n\to Z\to Y\)。纤维 \((n-1)\) 连通，故
	\[
		H_q(F_n)=0\ (0<q<n),\qquad H_n(F_n)\cong\pi_n(F_n)\cong\pi_n(X).
	\]
	在同调谱序列中，这使前 \(n-1\) 个正行为空。总次数 \(n\) 只可能有 \((n,0)\) 和 \((0,n)\) 两项；唯一可能击中后一项的箭头是
	\[
		\mathrm d_{n+1}:H_{n+1}(Y)=E^{n+1}_{n+1,0}\longrightarrow E^{n+1}_{0,n}=H_n(F_n).
	\]
	源端在前几页未改变，因为较早的去箭头都落入空行，零行又无来箭头。目标项的稳定余核是目标滤过的最浅子群，因而有正合列
	\begin{equation}\label{g2:eq:finite-generation-five-terms}
		H_{n+1}(Y)\xrightarrow{\mathrm d_{n+1}}H_n(F_n)
		\longrightarrow H_n(Z)\longrightarrow H_n(Y)\longrightarrow0.
	\end{equation}
	第一映射的像有限生成，因为 \(H_{n+1}(Y)\) 有限生成；它的余核是 \(H_n(Z)\cong H_n(X)\) 的子群，也有限生成。因此中间群 \(H_n(F_n)\cong\pi_n(X)\) 是两个有限生成群的扩张，仍有限生成，完成归纳。

	若 \(X\) 为有限 CW 复形，其胞腔链只有有限多个有限秩自由群，故所有同调群有限生成，最后的断言随即成立。
\end{proof}

证明的关键出现在式\eqref{g2:eq:finite-generation-five-terms}：未知的 \(\pi_n(X)\) 被放进一个纤维首个同调群，左边控制它被消去的部分，右边控制剩下的部分。两部分都有限生成，就足以控制整个未知群。这里无需算出传递的具体整数矩阵。

\ProseParagraphBreak

例如，对每个 \(m\ge2\) 和 \(n\ge2\)，\(\pi_n(S^m)\) 都是有限生成阿贝尔群。结论没有给出其阶，也没有断言它是有限群：\(\pi_m(S^m)=\mathbb Z\)，且前面已算出 \(\pi_3(S^2)=\mathbb Z\)。逐次有限生成与“所有高阶群都是有限群”是不同的断言。单连通假设也确实承担作用，章末将用 \(S^1\vee S^2\) 展示失去这一假设后的现象。

\subsection{Adams 与 Eilenberg–Moore 谱序列}
\label{g2:subsec:1105:03}

到这里，Serre 方法的输入与输出已经清楚：由基底和纤维的同调、运输作用及几何关系，求总空间的同调。路径纤维化使同一方法也能反求回路空间。若目标变成稳定同伦群，或已知基底与总空间而想系统地反求纤维的环结构，就需要不同的构造。

\ProseParagraphBreak

Adams\footnote{John Frank Adams（1930—1989），英国数学家，发展了计算稳定同伦群的谱序列方法，并在球面同伦与上同调运算方面作出重要贡献。}方法把模素数上同调与作用在它上面的自然运算一起使用。例如，两个分次向量空间相同的上同调，如果运算的作用不同，就会给出不同的代数输入。为此需要三项新结构：描述这些运算及其关系的代数；在该代数的模范畴中构造消解和高阶 Ext；把逐层消解实现为谱的塔，并证明它在相应素数完备化意义下收敛。第4章的 Ext 是理解消解的起点，但只计算整数模之间的一次 Ext，还没有包含这些结构。第13章将先讨论上同调运算，第9章的稳定化则提供稳定同伦的入口。

\ProseParagraphBreak

Eilenberg–Moore 方法面向同伦拉回。对于 \(F\to E\to B\)，纤维可以理解为 \(E\) 与一个点在 \(B\) 上的同伦拉回；相应的代数模型涉及把 \(C^*(E;k)\) 与 \(k\) 在 \(C^*(B;k)\) 上作导出张量积。这里 \(k\) 是域，“导出”表示先用适当的自由模型替换，以保留普通张量积会丢失的关系。它所需的新输入是带乘法的上链模型、条形消解，以及将该代数模型与实际同伦纤维比较的定理。在单连通与适当有限型条件下，第二页可用分次上同调代数上的 Tor 描述；这些比较与收敛条件须随新构造一同证明。

\ProseParagraphBreak

一个小的代数计算能说明为什么会出现这种输入。设 \(A=k[u]\)、\(|u|=2\)，增广把 \(u\) 送到零。\(k\) 作为分次 \(A\) 模有自由消解
\[
	0\longrightarrow A\varepsilon\xrightarrow{\ \varepsilon\mapsto u\ }A\longrightarrow k\longrightarrow0,
\]
其中 \(\varepsilon\) 的内部次数为二、消解次数为一。张量 \(k\) 后微分为零，所得 Tor 只有单位和一个总次数 \(2-1=1\) 的生成元；用 \(\varepsilon^2=0\) 的自由微分分次代数模型，可见其乘法为外代数。这个答案与 \(\Omega\mathbb CP^\infty\simeq S^1\) 的上同调吻合；后一个同伦型由 \(\mathbb CP^\infty=K(\mathbb Z,2)\)、路径长正合列及上面的模型唯一性得到。要对一般纤维化作同样的推断，承担几何与代数之间联系的正是新的比较定理。

\ProseParagraphBreak

本章的下一步则回到空间本身：有限截段 \(X[n]\to X[n-1]\) 已把一个同伦群分离为纤维，但仍未说明这些纤维怎样附着成原空间。第12章的障碍类与 Postnikov 分解将记录这部分数据，使这里用于证明有限性的构造进一步成为描述同伦类型的工具。
% 纲要条目排为编号小节标题，并进入目录与书签；说明段仍保留于纲要注释中。


\begin{chaptersummary}
	基底骨架的逆像给出 Serre 谱序列，第二页为基底取纤维同调或上同调局部系统后的群。沿路运输决定第二页是否能够使用常值系数，固定次数的骨架稳定及相容 CW 模型保证有限目标滤过。上同调页面的乘法满足总次数 Leibniz 规则，边缘分别对应投影拉回与纤维限制；传递则是纤维类延拓到下一基底骨架的障碍。Hopf 丛迫使相应传递为同构，定向球丛的两行给出 Euler 类与上同调正合列。可缩的道路空间迫使回路球面各次的微分为同构，进而决定加性群与整数除幂乘法。酉群的球面纤维序列由生成元次数迫使退化，再以自由群和乘积基解决环扩张。有限生成性经取子群、商和有限扩张保持；将单个同伦群放进 Eilenberg–MacLane 纤维后，这一性质由谱序列从同调传到全部高阶同伦群。
\end{chaptersummary}

\BookExercises

\begin{exercise}\label{g2:ex:serre-invariant-not-permanent}
	比较 \(S^1\to S^2\times S^1\to S^2\) 与 \(S^1\to S^3\to S^2\) 的上同调谱序列。证明两者所有纤维类都保持运输不变，但圆周的一次生成元只有在前一个纤维化中才是永久圈。解释“运输不变”与“能够从总空间限制下来”的区别。
\end{exercise}
\begin{solution}
	两者基底均单连通，所以运输作用平凡。第二页同为 \(\mathbb Z[x]/(x^2)\otimes\Lambda(u)\)，其中 \(x\) 位于 \((2,0)\)，\(u\) 位于 \((0,1)\)。唯一可能非零的微分为 \(\mathrm d_2u=kx\)。

	乘积中，第二投影拉回的圆周类限制到纤维恰为 \(u\)。边缘映射遂说明 \(u\) 是永久圈，\(k=0\)，乘积规则又使所有微分为零。Hopf 丛中，\(H^1(S^3)=H^2(S^3)=0\)，所以 \(\mathrm d_2\) 必须无核且无余核，即 \(k=\pm1\)。因此 \(u\) 在第二页被消去，不能是某个总空间上同调类的限制。运输不变只检查沿基底道路的相容性，使类能够进入第二页；更高骨架上的延拓还需要通过后续微分。
\end{solution}

\begin{exercise}\label{g2:ex:serre-torus-monodromy}
	令 \(f:T^2\to T^2\) 为 \(f([x],[y])=([-x],[-y])\)，记其映射环面为 \(M_f\)。用 Serre 谱序列求 \(H^k(M_f;\mathbb Z)\)，并说明哪一步需要处理加性扩张。
\end{exercise}
\begin{solution}
	\(H^0(T^2;\mathbb Z)=H^2(T^2;\mathbb Z)=\mathbb Z\)，\(H^1(T^2;\mathbb Z)=\mathbb Z^2\)。\(f^*\) 在一次是 \(-I\)，在二次是 \(\det(-I)=1\)，在零次也是恒等。因为 \(f^{-1}=f\)，上同调局部运输与这些作用相同。圆周的局部系数上链微分为作用的逆减恒等；因此第一行和第三行各在 \(p=0,1\) 有 \(\mathbb Z\)，中间一行的 \(p=0\) 项为零，\(p=1\) 项为 \((\mathbb Z/2)^2\)。只有两列，故 \(E_2=E_\infty\)。

	零、一、三次各只有一个非零分次块，分别给出 \(\mathbb Z\)。二次有短正合列
	\[
		0\longrightarrow(\mathbb Z/2)^2\longrightarrow H^2(M_f;\mathbb Z)\longrightarrow\mathbb Z\longrightarrow0.
	\]
	商群 \(\mathbb Z\) 自由，任选其生成元的一个提升就给出加性分裂，因而 \(H^2(M_f;\mathbb Z)\cong\mathbb Z\oplus(\mathbb Z/2)^2\)。其他次数为零。这里分裂由自由商保证；仅凭谱序列退化还不能跳过这一步。
\end{solution}

\begin{exercise}\label{g2:ex:serre-covering-coefficients}
	对 \(m\ge2\)，考虑覆盖 \(p:S^1\to S^1\)、\(p(z)=z^m\)。把它看成纤维有 \(m\) 个点的纤维化，写出第二页，并从局部系数映射解释 \(p^*:H^1(S^1;\mathbb Z)\to H^1(S^1;\mathbb Z)\) 为乘 \(m\)。
\end{exercise}
\begin{solution}
	纤维上同调只有零次非零，为 \(M=\mathbb Z^m\)。绕基底一周循环置换这 \(m\) 个点，因而在 \(M\) 上循环置换标准基，记该作用为 \(P\)；取相反循环置换也不改变以下核与余核。局部上链为 \(M\xrightarrow{P^{-1}-I}M\)，其核是对角子群 \(\mathbb Z(1,\ldots,1)\)，余核是所有标准基互相识别后得到的 \(\mathbb Z\)。因此 \(E_2^{0,0}\cong E_2^{1,0}\cong\mathbb Z\)，其余项为零，没有后续微分和加性扩张。

	投影拉回在系数上把一个常值整数 \(a\) 送到所有纤维点上的同一个值，即 \(a\mapsto(a,\ldots,a)\)。在一次局部上同调的余核中，\((1,\ldots,1)\) 是 \(m\) 个相同生成元之和，所以诱导映射为乘 \(m\)。这个算例同时说明，纤维不连通时，零行系数是 \(\mathcal H^0(F;\mathbb Z)\)，不能径直换成一份常值 \(\mathbb Z\)。
\end{solution}


\begin{exercise}\label{g2:ex:serre-global-fiber-basis}
	设 \(F\to E\xrightarrow{p}B\) 为道路连通空间的纤维化，\(B\) 为连通 CW 复形。假设 \(H^*(F;\mathbb Z)\) 有有限个齐次元素组成的整数基，且存在 \(e_1,\ldots,e_N\in H^*(E;\mathbb Z)\)，使它们在一个纤维上的限制组成这组基。证明
	\[
		\bigoplus_{j=1}^NH^{k-|e_j|}(B;\mathbb Z)\longrightarrow H^k(E;\mathbb Z),\qquad
		(b_j)_j\longmapsto\sum_jp^*b_j\smile e_j
	\]
	对每个 \(k\) 都是同构。说明结论为何依赖于这些实际存在的全局类。
\end{exercise}
\begin{solution}
	总空间上同调类在各纤维上的限制与道路运输相容，所以 \(e_j|_F\) 全部运输不变。它们构成每次纤维群的基，故整个局部系统为常值。有限秩自由系数的上链计算给出
	\[
		E_2^{p,q}=H^p(B;\mathbb Z)\otimes H^q(F;\mathbb Z).
	\]
	由边缘映射，每个 \(e_j|_F\) 都是永久圈。零行上的微分向负纤维次数出发，因而为零；每个第二页元素都是底类与这些纤维基元的线性组合，Leibniz 规则使全部第二微分为零。逐页重复，得到 \(E_2=E_\infty\)。

	给左侧按底类次数的递减滤过，右侧取 Serre 滤过。\(p^*b_j\) 的滤过至少为 \(|b_j|\)，而 \(e_j\) 的零列首项为 \(e_j|_F\)，所以题中映射保留滤过。它在每个伴随分次上正好是所列第二页同构。固定 \(k\)，滤过有限，对相邻滤过项的短正合列归纳就得到目标同构。这里也解决了加性扩张，因为选定的全局 \(e_j\) 已给出具体提升。若只有抽象的纤维基而没有全局提升，Hopf 丛的 \(u\) 就可能被微分消去，结论不成立。
\end{solution}

\begin{exercise}\label{g2:ex:circle-bundle-projective-plane-ring}
	设 \(S^1\to E\xrightarrow{p}\mathbb CP^2\) 为定向圆周丛，且 \(e(p)=k x\)，其中 \(x\) 生成 \(H^2(\mathbb CP^2;\mathbb Z)\)、\(k\ne0\)。求全部整数上同调群及杯积环。特别指出 \(|k|>1\) 时二次挠类的平方。
\end{exercise}
\begin{solution}
	第二页为 \(\mathbb Z[x]/(x^3)\otimes\Lambda(u)\)，其中 \(|x|=(2,0)\)、\(|u|=(0,1)\)。按 Euler 类定义，\(\mathrm d_2u=kx\)，所以
	\[
		\mathrm d_2(xu)=kx^2,\qquad \mathrm d_2(x^2u)=0.
	\]
	前两张微分在整数群上单射，分别留下 \(\mathbb Z/|k|\) 的余核。第三页以后没有可能的微分，因此总上同调为
	\[
		H^j(E;\mathbb Z)\cong
		\begin{cases}
			\mathbb Z,&j=0,5,\\
			\mathbb Z/|k|,&j=2,4,\\
			0,&\text{其余 }j.
		\end{cases}
	\]
	每个次数只有一层非零滤过。令 \(a=p^*x\)，再选 \(v\in H^5(E)\) 对应 \(x^2u\) 的生成元。乘法与伴随分次相容，故 \(a^2\) 生成四次群；当 \(|k|>1\) 时尤其非零。其余正次乘积落在五次以上或含 \(a^3\)，均为零。完整环可写成以 \(|a|=2,|v|=5\) 为生成元的分次交换环，关系为 \(ka=0,a^3=0,av=0,v^2=0\)。\(k=\pm1\) 时 \(a=0\)，只留下零次与五次整数群。
\end{solution}

\begin{exercise}\label{g2:ex:serre-hopf-pullback-degree}
	令 \(p:S^3\to S^2\) 为复 Hopf 丛，选择生成元 \(u\in H^2(S^2;\mathbb Z)\) 使 \(e(p)=u\)。对整数 \(k\ne0\)，取度为 \(k\) 的映射 \(f_k:S^2\to S^2\)，记拉回圆周丛为 \(p_k:E_k\to S^2\)。求 \(E_k\) 的全部整数上同调群与正次杯积，并证明 \(p_k\) 没有截面。再比较 \(k=0\) 时的加性群。
\end{exercise}
\begin{solution}
	自然性给出 \(e(p_k)=f_k^*u=ku\)。球丛正合列中的唯一非零乘 Euler 类映射为 \(H^0(S^2)=\mathbb Z\xrightarrow{\times k}H^2(S^2)=\mathbb Z\)。因此当 \(k\ne0\) 时，
	\[
		H^j(E_k;\mathbb Z)=
		\begin{cases}
			\mathbb Z,&j=0,3,\\
			\mathbb Z/|k|\mathbb Z,&j=2,\\
			0,&\text{其余 }j.
		\end{cases}
	\]
	所有正次类的乘积都落在至少四次，故全为零。由于 \(e(p_k)=ku\ne0\)，截面消失结论排除了连续截面。

	当 \(k=0\) 时，正合列中的乘 Euler 类映射全为零，给出 \(H^j(E_0;\mathbb Z)=\mathbb Z\)（\(j=0,1,2,3\)），其余为零。每个次数至多有一层非零滤过，故这里没有加性延拓问题。该结果与 \(S^2\times S^1\) 的加性上同调相同；也可取常值 \(f_0\)，其拉回直接就是这个乘积丛。
\end{solution}

\begin{exercise}\label{g2:ex:even-sphere-euler-two-torsion}
	设 \(m\ge2\) 为偶数，\(p:E\to B\) 为道路连通 CW 复形上的定向 \(S^m\) 丛。证明 \(2e(p)=0\)。由此说明：若 \(H^{m+1}(B;\mathbb Z)\) 没有 \(2\) 挠，则 \(e(p)=0\)；特别地，定向 \(S^2\) 丛在底空间 \(S^3\) 上不可能具有非零 Euler 类。
\end{exercise}
\begin{solution}
	令 \(u\in E_2^{0,m}\) 为纤维定向生成元。因为纤维没有 \(2m\) 次上同调，\(u^2=0\)。在第 \(m+1\) 页以前没有非零微分，故该等式在 \(E_{m+1}\) 仍成立。\(m\) 为偶数，Leibniz 规则以及 \(u\) 与 \(e(p)\) 的交换性给出
	\[
		0=\mathrm d_{m+1}(u^2)=e(p)u+u e(p)=2e(p)u.
	\]
	乘 \(u\) 把 \(H^{m+1}(B;\mathbb Z)\) 同构地识别为顶行 \(E_{m+1}^{m+1,m}\)，所以 \(2e(p)=0\)。无 \(2\) 挠时这就迫使 \(e(p)=0\)。对于 \(B=S^3\)、\(m=2\)，所需群为 \(H^3(S^3;\mathbb Z)=\mathbb Z\)，故结论适用。这里不能把 \(m\) 改为奇数：奇数时 Leibniz 规则的两项互相抵消，而 Hopf 丛的 Euler 类可以有无穷阶。
\end{solution}

\begin{exercise}\label{g2:ex:divided-power-finite-generation}
	设 \(m\ge2\) 为偶数，记 \(A=H^*(\Omega S^{m+1};\mathbb Z)=\Gamma_{\mathbb Z}(a_1)\)。证明 \(A\) 在每个次数都是有限生成阿贝尔群，但作为 \(\mathbb Z\) 代数不是有限生成的。再证明换成有理系数后，\(H^*(\Omega S^{m+1};\mathbb Q)\) 是一个元素生成的多项式代数。
\end{exercise}
\begin{solution}
	每个次数的群为零或 \(\mathbb Z\)，第一项已成立。若 \(A\) 由有限多个元素作为代数生成，把这些元素写成加性基 \(a_j\) 的有限和，取 \(N\) 使其中出现的下标都不超过 \(N\)。它们生成的子代数包含在由 \(a_0,a_1,\ldots,a_N\) 生成的子代数中。选素数 \(\ell>N\)，再把系数约化到 \(\mathbb F_\ell\)。

	任何由 \(a_1,\ldots,a_N\) 构成、次数为 \(\ell m\) 的单项式，其下标之和为 \(\ell\)，各正下标都严格小于 \(\ell\)。重复应用除幂乘法，它等于
	\[
		a_{j_1}\cdots a_{j_s}=\frac{\ell!}{j_1!\cdots j_s!}\,a_\ell,
		\qquad j_1+\cdots+j_s=\ell,\quad 1\le j_i\le N.
	\]
	分子含素因子 \(\ell\)，每个分母都不含，故整数系数被 \(\ell\) 整除。于是所有这些单项式在 \(A\otimes\mathbb F_\ell\) 中为零，而 \(a_\ell\) 的像非零，矛盾。非齐次候选生成元也没有例外，因为取它们的齐次分量只会扩大生成的子代数。

	整数同调群均自由，系数定理把有理上同调识别为 \(A\otimes\mathbb Q\)，且系数映射保持杯积。等式 \(a_1^j=j!a_j\) 此时允许除以 \(j!\)，说明 \(a_1\) 生成整个环。各 \(a_1^j\) 位于不同次数且非零，没有非零多项式关系，故得到 \(\mathbb Q[a_1]\)。这区分了逐次数的有限生成性与整个分次代数的有限生成性。
\end{solution}


\begin{exercise}\label{g2:ex:cohomology-suspension-products}
	设 \(B\) 为带基点的道路连通 CW 复形，\(\alpha\in\widetilde H^a(B;\mathbb Z)\)、\(\beta\in\widetilde H^b(B;\mathbb Z)\)，其中 \(a,b>0\)。证明上同调悬挂满足 \(\sigma(\alpha\smile\beta)=0\)。用 \(B=\mathbb CP^\infty\) 说明它为何不能在所有次数上成为同构。
\end{exercise}
\begin{solution}
	先说明悬挂空间的任意两个正次类的乘积为零。将悬挂 \(\Sigma Y\) 用稍加延长的上下两个开锥 \(U,V\) 覆盖，两开集各自可缩。任意正次类 \(c,d\) 分别在 \(U,V\) 上限制为零，所以由空间对正合列可选相对提升
	\[
		\widetilde c\in H^*(\Sigma Y,U),\qquad \widetilde d\in H^*(\Sigma Y,V).
	\]
	相对杯积落在 \(H^*(\Sigma Y,U\cup V)=H^*(\Sigma Y,\Sigma Y)=0\)，故绝对乘积 \(cd=0\)。

	对 \(Y=\Omega B\)，求值给出 \(\operatorname{ev}:\Sigma\Omega B\to B\)。杯积自然性使
	\[
		\operatorname{ev}^*(\alpha\beta)=\operatorname{ev}^*\alpha\smile\operatorname{ev}^*\beta=0.
	\]
	再作悬挂次数同构，按正文对 \(\sigma\) 的识别得到 \(\sigma(\alpha\beta)=0\)。在 \(B=\mathbb CP^\infty\) 时，二次生成元 \(u\) 的首次悬挂生成 \(H^1(\Omega B)=\mathbb Z\)，但 \(u^j\ne0\)（\(j\ge2\)）全是正次类的乘积，故 \(\sigma(u^j)=0\)。因此首次非零次数的同构不能不加条件地延伸到高次。
\end{solution}

\begin{exercise}\label{g2:ex:unitary-determinant-product}
	设 \(n\ge2\)。证明
	\[
		\Phi:S^1\times\operatorname{SU}(n)\longrightarrow\operatorname{U}(n),\qquad(z,A)\longmapsto\operatorname{diag}(z,1,\ldots,1)A
	\]
	为同胚，写出其逆，并说明行列式在这个乘积中对应哪张投影。这一同胚一般是否为直积群到酉群的群同态？比较所得上同调环与\cref{g2:thm:unitary-cohomology}。
\end{exercise}
\begin{solution}
	记 \(D(z)=\operatorname{diag}(z,1,\ldots,1)\)。逆映射为
	\[
		U\longmapsto\bigl(\det U,D(\det U)^{-1}U\bigr).
	\]
	第二分量的行列式为一，两张公式连续且互逆，因此 \(\Phi\) 为同胚；\(\det\Phi(z,A)=z\)，所以行列式对应第一投影。在直积群中 \((z,I)\) 与 \((1,A)\) 总交换，但它们的像 \(D(z)\) 与 \(A\) 未必交换。例如在前两个坐标中取 \(A=\left(\begin{smallmatrix}0&1\\-1&0\end{smallmatrix}\right)\)，其他坐标为单位阵；当 \(z\ne1\) 时，两张矩阵不交换。因此 \(\Phi\) 一般不是群同态。

	两因子的整数上同调都自由，Künneth 定理及外积乘法给出
	\[
		H^*(\operatorname{U}(n);\mathbb Z)\cong H^*(S^1;\mathbb Z)\otimes H^*(\operatorname{SU}(n);\mathbb Z)
		\cong\Lambda_{\mathbb Z}(x_1,y_3,\ldots,y_{2n-1}).
	\]
	其中一次类正是行列式拉回的生成元，与正文外代数一致。这个结论比较的是空间的杯积环，并未要求所选同胚保持群乘法。
\end{solution}

\begin{exercise}\label{g2:ex:unitary-four-middle-pairing}
	在 \(\operatorname{U}(4)\) 中选择\cref{g2:thm:unitary-cohomology}的整数生成元，并将 \(v=x_1x_3x_5x_7\) 取作十六次生成元。求 \(H^8(\operatorname{U}(4);\mathbb Z)\) 的一组基，并写出杯积 \(H^8\times H^8\rightarrow H^{16}=\mathbb Zv\) 对应的整数矩阵。
\end{exercise}
\begin{solution}
	在 \(1,3,5,7\) 中选取互异项，和等于 \(8\) 的方式只有 \(1+7\) 与 \(3+5\)。所以一组基为 \(a=x_1x_7\)、\(b=x_3x_5\)。平方 \(a^2,b^2\) 都含某个奇次生成元的平方，故均为零。而
	\[
		ab=x_1x_7x_3x_5=(-1)^2x_1x_3x_5x_7=v.
	\]
	两次交换把 \(x_7\) 移到末尾，因而没有负号。\(|a|=|b|=8\) 还给出 \(ba=ab\)，所以矩阵为 \(\left(\begin{smallmatrix}0&1\\1&0\end{smallmatrix}\right)\)。其行列式为 \(-1\)，与闭定向十六维流形中维整数交形式的单模性相符；若把顶次生成元换成 \(-v\)，矩阵整体变号。
\end{solution}


\begin{exercise}\label{g2:ex:hurewicz-mod-finite-generation}
	设 \(X\) 为单连通 CW 复形，\(n\ge2\)，且 \(\pi_i(X)\) 对 \(2\le i<n\) 都有限生成。证明 Hurewicz 同态 \(h:\pi_n(X)\to H_n(X;\mathbb Z)\) 的核与余核都有限生成。这里不预先假定 \(\pi_n(X)\) 或 \(H_n(X)\) 有限生成。
\end{exercise}
\begin{solution}
	\(n=2\) 时 Hurewicz 同态为同构，核与余核为零。设 \(n>2\)，按正文构造 \(F_n\to X[n]\to X[n-1]\)，记总空间为 \(Z\)、基底为 \(Y\)。低阶同伦群有限生成使 \(Y\) 的全部同调逐次有限生成。\(F_n\) 为 \((n-1)\) 连通，且 \(\pi_n(F_n)\to\pi_n(Z)\) 为同构。Hurewicz 自然性给出交换图
	\[
		\begin{tikzcd}
			\pi_n(F_n)\arrow[r,"\cong"]\arrow[d,"\cong"']&\pi_n(Z)\arrow[d,"h_Z"]\\
			H_n(F_n)\arrow[r]&H_n(Z).
		\end{tikzcd}
	\]
	下箭头的核由\cref{g2:eq:finite-generation-five-terms}是 \(H_{n+1}(Y)\) 的像，余核同构于 \(H_n(Y)\)，两者均有限生成。因此右箭头也有有限生成的核与余核。包含 \(X\to X[n]\) 在 \(\pi_n\) 和 \(H_n\) 上都是同构，并与 Hurewicz 映射交换，所以同一结论适用于 \(h\)。这正是 Hurewicz 映射模 \(\mathcal C_{\mathrm{fg}}\) 为同构的含义。
\end{solution}

\begin{exercise}\label{g2:ex:finite-homology-infinite-pi-two}
	令 \(X=S^1\vee S^2\)，基点取楔和点。构造其万有覆盖，并证明 \(H_k(X;\mathbb Z)\) 对每个 \(k\) 都有限生成，但
	\[
		\pi_2(X)\cong\bigoplus_{j\in\mathbb Z}\mathbb Z.
	\]
	说明这个例子对有限生成定理的单连通条件有何意义。
\end{exercise}
\begin{solution}
	把实直线按整数点分成一胞腔，并在每个整数点 \(j\) 附着一个球面 \(S_j^2\)。所得空间记为 \(\widetilde X\)。将实直线按 \(t\mapsto e^{2\pi\mathrm it}\) 投到 \(S^1\)，每个 \(S_j^2\) 同胚地投到 \(S^2\)，得到覆盖 \(\widetilde X\rightarrow X\)：在楔和点附近，圆周的一段小弧连同球面的一小片，其原像是围绕各整数点的互不相交同胚副本。对其余点，覆盖性质直接由圆周覆盖及球面同胚得到。

	\(\widetilde X\) 的一维骨架为实直线，而每个二胞腔沿常值映射附着，因此其基本群为零，所构造的覆盖是万有覆盖。胞腔链中没有三胞腔，二胞腔边界也为零，所以
	\[
		H_2(\widetilde X;\mathbb Z)=\bigoplus_{j\in\mathbb Z}\mathbb Z[S_j^2].
	\]
	\cref{g2:thm:hurewicz-absolute}用于单连通的 \(\widetilde X\)，得到 \(\pi_2(\widetilde X)\cong H_2(\widetilde X;\mathbb Z)\)。覆盖映射的离散纤维在正阶同伦群中为零，纤维化长正合列因而给出 \(\pi_2(\widetilde X)\cong\pi_2(X)\)，得到题中同构。这个群不是有限生成的，因为有限多个元素只涉及有限多个直和坐标。

	另一方面，\(X\) 有一个零胞腔、一个一胞腔和一个二胞腔，胞腔边界全为零。因此 \(H_0(X)=H_1(X)=H_2(X)=\mathbb Z\)，更高同调为零。\(X\) 虽然本身是有限 CW 复形，\(\pi_2(X)\) 仍不是有限生成阿贝尔群；单连通条件不能在有限生成定理中直接删去。绕圆周一圈把 \(S_j^2\) 移到 \(S_{j+1}^2\)，正显示了基本群作用如何产生无穷多个整数生成元。
\end{solution}

\begin{exercise}\label{g2:ex:finite-fundamental-group-finiteness}
	设 \(X\) 为连通有限 CW 复形，且 \(\pi_1(X)\) 为有限群。证明所有 \(\pi_n(X)\)（\(n\ge2\)）都是有限生成阿贝尔群。
\end{exercise}
\begin{solution}
	令 \(m=|\pi_1(X)|\)。万有覆盖 \(\widetilde X\rightarrow X\) 有 \(m\) 层。将 \(X\) 的胞腔特征映射提升到覆盖空间，得到 \(\widetilde X\) 的 CW 结构；每个底胞腔恰有 \(m\) 个提升胞腔。因此 \(\widetilde X\) 也是有限 CW 复形，其每个整数同调群都有限生成。它单连通，\cref{g2:thm:serre-finite-generation}于是说明 \(\pi_n(\widetilde X)\) 对所有 \(n\ge2\) 都有限生成。离散纤维的长正合列给出 \(\pi_n(\widetilde X)\cong\pi_n(X)\)，完成证明。
\end{solution}

\begin{exercise}\label{g2:ex:finite-type-not-finite-cw}
	设 \(m>1\) 为整数，\(G=\mathbb Z/m\)。利用本节的有限群模型说明 \(K(G,1)\) 的每个整数同调群都有限生成，再证明它不可能具有有限 CW 同伦型。
\end{exercise}
\begin{solution}
	本节构造的 \(BG\) 每维只有有限多个胞腔，因为该维单形由有限群 \(G\) 的有限长度序列决定。因此每个胞腔链群都是有限秩自由阿贝尔群，圈群为其子群，仍有限生成，取商后每个 \(H_k(BG;\mathbb Z)\) 都有限生成。\(BG\) 是 \(K(G,1)\)，同伦不变性便给出前一结论。

	假设某个连通有限 CW 复形 \(Y\) 与 \(K(G,1)\) 同伦等价。其万有覆盖 \(\widetilde Y\) 有 \(m\) 层，每个胞腔提升成 \(m\) 个胞腔，所以 \(\widetilde Y\) 也有限，且胞腔计数给出
	\[
		\chi(\widetilde Y)=m\chi(Y).
	\]
	\(Y\) 的全部高阶同伦群为零，覆盖映射在这些群上诱导同构；又因 \(\widetilde Y\) 单连通，它的全部正阶同伦群都为零。\cref{g2:thm:whitehead}因而说明 \(\widetilde Y\) 可缩，所以 \(\chi(\widetilde Y)=1\)。于是 \(1=m\chi(Y)\)，与 \(m>1\) 及 \(\chi(Y)\in\mathbb Z\) 矛盾。逐次同调有限生成允许非零信息出现于无穷多个次数，也并不保证存在只有有限多个胞腔的模型。
\end{solution}

